%%%PAGE 68 rot=0 legibility=good summary=曲线运动章节提纲(思维导图式)：曲线运动条件与a、v关系，运动的合成与分解(小船渡河、关联速度图)，抛体运动各类问题，圆周运动与传动方式
%%%BLOCK id=068-01 ch=04 sec=曲线运动·条件与速度变化 kind=knowledge alt=none cont=none y=0.07-0.25
% 原稿：页面最上方一行（位于“Date. NO.”表头上方），以及表头下方右侧两行、其下两处批注
$a$\ $V$锐角\ $V\uparrow$\qquad $a$\ $V$钝角\ $V\downarrow$\qquad $a$\ $V$\unclear{垂直}\ 速度大小不变，方向改变 % CHECK: 顶行三项中 a 与 V 之间似无符号（应为夹角关系），第三项“垂直”二字不确定

$a\perp V$\quad 改变$V$方向\quad 法向加速度

$a\parallel V$\quad 改变$V$大小\quad 切向加速度

$a$、$V$不同方向

力与速度夹轨迹\quad 凹侧方向偏向力
%%%END
%%%BLOCK id=068-02 ch=04 sec=章节提纲 kind=summary alt=none cont=none y=0.12-0.97
% 版面：左侧一列为各部分标题（曲线运动；运动的合成与分解，其下一行小字“合成、分解”；抛体运动，其下“分解”；圆周运动，其下一行小字），右侧为对应条目
\begin{itemize}
  \item \textbf{曲线运动}
  \begin{itemize}
    \item 曲线运动、骑马射箭、转台投篮、风中骑车、\unclear{滑}板运动
  \end{itemize}
  \item \textbf{运动的合成与分解}\\ 合成、分解
  \begin{itemize}
    \item 运动的合成与分解，\quad 等时性\quad 独立性\quad 等效性
    \item 小船渡河模型、
    \item 关联速度模型、 % 图及批注见 068-03
  \end{itemize}
  \item \textbf{抛体运动}\\ 分解
  \begin{itemize}
    \item （类）平抛运动
    \item （类）斜抛运动
    \item 落点有约束的平抛运动
    \item 平抛运动的临界问题
    \item 斜抛运动问题。
  \end{itemize}
  \item \textbf{圆周运动}\\ \unclear{受力}
  \begin{itemize}
    \item 线速度\quad \underline{角速度}\quad \underline{周期}、\underline{转速}\quad 向心加速度\quad 向心力
    \item 离心运动（$F_n<m\omega^2 r$）、和向心运动（$F_n>m\omega^2 r$） % 原稿两式分别写在“离心运动”“和向心运动”上方
    \item $\underbrace{\text{转杆}}_{\text{同}\,\omega}$、$\underbrace{\text{皮带传动、齿轮传动、摩擦传动}}_{\text{同}\,v}$、$\underset{\text{同}\,\omega}{\text{同轴传动}}$ % 原稿批注写在各项下方：转杆下波浪线+“同ω”；皮带/齿轮/摩擦传动下长波浪线+“同v”；同轴传动下“同ω”
  \end{itemize}
\end{itemize}
%%%END
%%%BLOCK id=068-03 ch=04 sec=关联速度模型 kind=diagram alt=none cont=none y=0.34-0.46
\textbf{关联速度模型}

%FIG p068_a 0.335 0.346 0.494 0.415
\notefig[0.28]{figures/p068_a.png}

%FIG p068_b 0.49 0.366 0.654 0.45
\notefig[0.28]{figures/p068_b.png}

沿绳垂绳，沿杆垂杆

力和速度的分解不同
%%%END
%%%PAGEEND 68
